(1) 取自然单位 $\hbar=c=1$，真空混合角 $\theta$、质量平方差 $\Delta m^2>0$、中微子能量 $E$，定义 $\Delta=\Delta m^2/(2E)$。均匀物质电子势 $V$。写出味基底去迹的二能级哈密顿量。
(2) 求物质有效混合角与本征能级差 $\Omega$。
(3) 求初始电子中微子传播长度 $L$ 后转化为另一味的概率。
(4) 给出共振密度条件，讨论真空和反中微子极限。
